\documentclass[10pt,a4paper]{amsart}
\usepackage[margin=19mm]{geometry}
\usepackage{amsmath,amssymb,mathtools}
\usepackage[scr=rsfs,cal=euler]{mathalfa}
\usepackage{xfrac}
\usepackage[unicode,hidelinks,bookmarks=false]{hyperref}
\newcommand{\ie}{i.e.\ }
\newcommand{\cf}{cf.\ }
\newcommand{\resp}{respectively\ }
\newcommand{\Subsection}{\subsection}
\providecommand{\nobreakdash}{\nobreak\hbox{-}}
\setlength{\emergencystretch}{2em}
\multlinegap0pt
\usepackage{xcolor}
\usepackage{amssymb}
\usepackage{tikz}
\newtheorem{theorem}{Theorem}
\newtheorem{corollary}{Corollary}
\title{Why Self-Supervised Encoders Want to Be Normal}
\author{Yuval Domb}
\date{Source: arXiv:2604.27743v2 (CC BY 4.0)}
\begin{document}
\maketitle

\noindent\emph{Source and licence.} Why Self-Supervised Encoders Want to Be Normal. Version arXiv:2604.27743v2 (CC BY 4.0), distributed by arXiv under the Creative Commons Attribution 4.0 International licence, http://creativecommons.org/licenses/by/4.0/, as recorded in the arXiv licence metadata for this exact version and shown on the abstract page https://arxiv.org/abs/2604.27743v2. Full licence notice: \texttt{licenses/arxiv-2604.27743-CC-BY-4.0.txt}.\par\smallskip

\noindent\textit{Editorial scope.} Counted unit: the article's corollary cor:min\_rate (source0.tex lines 708--713). The article's complete original proof of it is delivered with it (source0.tex lines 715--723, one \texttt{proof} environment of 9 lines). No mathematics is written, repaired or re-derived: the statement and the proof are the article's own lines, in the article's own notation, and no retained line is substituted or re-flowed.  Retained uncounted as the source-local context the proof needs: lines 637--642; lines 669--691.  Nothing was omitted from any retained span, so the package carries no omission record.  No retained line contains a citation, so the wrapper carries no bibliography and no bibliography entry.  No retained line contains a figure environment, an image-inclusion command or drawing code, so no image is delivered and the package declares no asset dependency.  Editorial formatting only, in addition: an AMS \texttt{amsart} preamble that restates the article's own declarations and macros used by the retained lines, namely the environments \texttt{theorem}, \texttt{corollary} on separate counters, together with the standard packages those lines need.  The retained environments print in this document as Theorem 1, Corollary 1 in order of appearance, whereas the article prints the same items with the numbering of its own full document.  The complete native source of the article as distributed in the arXiv e-print for this exact version is bundled unchanged as \texttt{sources/199496/source0.tex}.
\begin{align}
    p(w) &= \sum_x p(x)\,p(w|x),
    \label{eq:marginal}\\
    p(y|w) &= \frac{1}{p(w)}\sum_x p(x)\,p(w|x)\,p(y|x).
    \label{eq:decoder}
\end{align}

\begin{theorem}[Minimal sufficient statistic]
\label{thm:mss}
Define $f^*:\mathcal{X}\to\mathcal{P}_K$ by $f^*(x):=p(Y|x)$, and let
$\{C_i\}$ denote the partition of $\mathcal{X}$ into level sets of $f^*$:
\begin{align}
    x \sim x' \quad \Longleftrightarrow \quad f^*(x)=f^*(x'),
    \label{eq:equiv}
\end{align}
where $C_i=\{x:f^*(x)=q_i\}$ for some $q_i\in\mathcal{P}_K$.
Define the random variable
\begin{align}
    W^* := f^*(X),
    \label{eq:mss}
\end{align}
so that $W^*$ takes values in $\mathcal{P}_K$. Then:
\begin{enumerate}
    \item $W^*$ is a sufficient statistic of $X$ for $Y$:
    $I(X;Y|W^*)=0$.
    \item $W^*$ is minimal: any encoder $W$ with $I(X;Y|W)=0$
    satisfies $W^*=g(W)$ almost surely, where $g(w):=p(Y|w)$ is the
    self-consistent decoder \eqref{eq:decoder}.
\end{enumerate}
\end{theorem}

\begin{corollary}[Minimum rate]
\label{cor:min_rate}
$W^*$ achieves the minimum rate among all sufficient encoders:
$I(X;W^*)=H(W^*)$, and any $W$ with $I(X;Y|W)=0$ satisfies
$I(X;W)\geq H(W^*)$, with equality if and only if $I(X;W|W^*)=0$.
\end{corollary}

\begin{proof}
By Theorem~\ref{thm:mss}, $W^*=g(W)$, so the data processing
inequality applied to $X\to W\to W^*$ gives
$I(X;W)\geq I(X;W^*)$.
Since $W^*=f^*(X)$ is deterministic, $H(W^*|X)=0$, hence
$I(X;W^*)=H(W^*)$.
Equality $I(X;W)=H(W^*)$ holds if and only if $I(X;W|W^*)=0$,
by the chain rule $I(X;W)=H(W^*)+I(X;W|W^*)$.
\end{proof}

\end{document}
